% GENERATED FILE. Edit canonical lessons, then run npm run generate:books.
% canonical-source-hash: fnv1a64:e368dc0c2eb60186
% canonical-lessons: JA-C42-kawa, JA-C42-umi, JA-C42-kusa, JA-C42-mori, JA-C42-ishi

\chapter{River, Sea, Grass, Forest, Stone}
\label{ch:ja-river-sea-grass-forest-stone}
\hlchaptermodality{42}

\begin{chapteropening}
\textbf{By the end of this chapter:} \emph{I can say a river, the sea, grass, a forest and a stone.}

Complete the last lesson of chapter 42: i can say a river, the sea, grass, a forest and a stone.
\end{chapteropening}

\section[kawa]{\ja{かわ} (kawa) — a river}

\label{lesson:JA-C42-kawa}



\begin{quote}
Before the new one: say the Japanese for money, then the Japanese for a mountain.
\end{quote}

\subsection*{You'll want to know: \ja{かわ}}
\textbf{\ja{かわ}} — \emph{kawa} — \textquotedblleft{}a river\textquotedblright{}.

\begin{grammarlens}[title={What you've built}]
One more word for this chapter.
\end{grammarlens}

\subsection*{Your turn}
\begin{itemize}
\raggedright
  \item \emph{Say it:} \emph{kawa}
  \item \emph{Say it:} \emph{kawa}, once more
  \item \emph{Say it:} say \emph{kawa}
  \item \emph{Recall:} say the Japanese for money, then the Japanese for a mountain, then say \emph{kawa} again
\end{itemize}

\begin{tcolorbox}[breakable,colback=teal!4,colframe=teal!35!black,title={Before you move on}]
What does \ja{かわ} mean? (\textquotedblleft{}A river\textquotedblright{}.) Say it once more.
\end{tcolorbox}
\section[umi]{\ja{うみ} (umi) — the sea}

\label{lesson:JA-C42-umi}



\begin{quote}
Before the new one: say the Japanese for a mountain, then the Japanese for a river.
\end{quote}

\subsection*{You'll want to know: \ja{うみ}}
\textbf{\ja{うみ}} — \emph{umi} — \textquotedblleft{}the sea\textquotedblright{}.

\begin{grammarlens}[title={What you've built}]
One more word for this chapter.
\end{grammarlens}

\subsection*{Your turn}
\begin{itemize}
\raggedright
  \item \emph{Say it:} \emph{umi}
  \item \emph{Say it:} \emph{umi}, once more
  \item \emph{Say it:} say \emph{umi}
  \item \emph{Recall:} say the Japanese for a mountain, then the Japanese for a river, then say \emph{umi} again
\end{itemize}

\begin{tcolorbox}[breakable,colback=teal!4,colframe=teal!35!black,title={Before you move on}]
What does \ja{うみ} mean? (\textquotedblleft{}The sea\textquotedblright{}.) Say it once more.
\end{tcolorbox}
\section[kusa]{\ja{くさ} (kusa) — grass}

\label{lesson:JA-C42-kusa}



\begin{quote}
Before the new one: say the Japanese for a river, then the Japanese for the sea.
\end{quote}

\subsection*{You'll want to know: \ja{くさ}}
\textbf{\ja{くさ}} — \emph{kusa} — \textquotedblleft{}grass\textquotedblright{}.

\begin{grammarlens}[title={What you've built}]
One more word for this chapter.
\end{grammarlens}

\subsection*{Your turn}
\begin{itemize}
\raggedright
  \item \emph{Say it:} \emph{kusa}
  \item \emph{Say it:} \emph{kusa}, once more
  \item \emph{Say it:} say \emph{kusa}
  \item \emph{Recall:} say the Japanese for a river, then the Japanese for the sea, then say \emph{kusa} again
\end{itemize}

\begin{tcolorbox}[breakable,colback=teal!4,colframe=teal!35!black,title={Before you move on}]
What does \ja{くさ} mean? (\textquotedblleft{}Grass\textquotedblright{}.) Say it once more.
\end{tcolorbox}
\section[mori]{\ja{もり} (mori) — a forest}

\label{lesson:JA-C42-mori}



\begin{quote}
Before the new one: say the Japanese for the sea, then the Japanese for grass.
\end{quote}

\subsection*{You'll want to know: \ja{もり}}
\textbf{\ja{もり}} — \emph{mori} — \textquotedblleft{}a forest\textquotedblright{}.

\begin{grammarlens}[title={What you've built}]
One more word for this chapter.
\end{grammarlens}

\subsection*{Your turn}
\begin{itemize}
\raggedright
  \item \emph{Say it:} \emph{mori}
  \item \emph{Say it:} \emph{mori}, once more
  \item \emph{Say it:} say \emph{mori}
  \item \emph{Recall:} say the Japanese for the sea, then the Japanese for grass, then say \emph{mori} again
\end{itemize}

\begin{tcolorbox}[breakable,colback=teal!4,colframe=teal!35!black,title={Before you move on}]
What does \ja{もり} mean? (\textquotedblleft{}A forest\textquotedblright{}.) Say it once more.
\end{tcolorbox}
\section[ishi]{\ja{いし} (ishi) — a stone}

\label{lesson:JA-C42-ishi}



\begin{quote}
Before the new one: say the Japanese for grass, then the Japanese for a forest.
\end{quote}

\subsection*{You'll want to know: \ja{いし}}
\textbf{\ja{いし}} — \emph{ishi} — \textquotedblleft{}a stone\textquotedblright{}.

\begin{grammarlens}[title={What you've built}]
One more word for this chapter.
\end{grammarlens}

\subsection*{Your turn}
\begin{itemize}
\raggedright
  \item \emph{Say it:} \emph{ishi}
  \item \emph{Say it:} \emph{ishi}, once more
  \item \emph{Say it:} say \emph{ishi}
  \item \emph{Recall:} say the Japanese for grass, then the Japanese for a forest, then say \emph{ishi} again
\end{itemize}

\begin{tcolorbox}[breakable,colback=teal!4,colframe=teal!35!black,title={Before you move on}]
What does \ja{いし} mean? (\textquotedblleft{}A stone\textquotedblright{}.) Say it once more.
\end{tcolorbox}
